% command to make sure colorboxes don't create whitespace between code lines
\fboxsep=.6pt

% default settings for listings
\lstset{style=python}

\title[Programmieren]{Programmieren: Datenstrukturen}

\subtitle{Listen}

\author[Wendl, Cyril] % (optional)
{Cyril Wendl}

\institute[KLW] % (optional)
{
	Fachschaft Informatik\\
	Kantonsschule im Lee
}

\date[\today] % (optional)

\logo{\includegraphics[width=2cm]{\thelogo}}

\titlegraphic{\vspace{-.7cm}\includegraphics[width=\textwidth]{"GrundlagenInfo/06_Datenbanken/Figures/title-emojis"}}

\frame{\titlepage}

\begin{frame}[fragile]{Beispiel}
	\begin{columns}
		\begin{column}{.5\textwidth}
			\begin{lstlisting}[basicstyle=\ttfamily\scriptsize]
# zwei Zahlen summieren
def summiere(x1, x2):
	summe = x1 + x2
	print(summe)
	
summiere(5, 16)
\end{lstlisting}
		\end{column}

		\begin{column}{.5\textwidth}
			\begin{lstlisting}[basicstyle=\ttfamily\scriptsize]
# viele Zahlen summieren
def summiere(x1, x2, x3, x4, ...):
	summe = x1 + x2 + x3 + x4 + ...
	print(summe)
	
summiere(3, 18, 2, 4, ...)?
\end{lstlisting}
		\end{column}
	\end{columns}
	\begin{itemize}[<+->]
		\item Wie kann man \lstinline|summiere(x1, x2, …, x100)| umsetzen? $\rightarrow$ umständlich!
		\item \lstinline|summiere(...)| für eine beliebige Anzahl von Zahlen? $\rightarrow$ unmöglich mit aktuellem Wissen!
		\item Lösung: \textbf{Listen}
	\end{itemize}
\end{frame}




\begin{frame}[fragile]{Listen}{Listen erstellen und abfragen}
	\begin{lstlisting}
# Neue Variable des Typs Liste erstellen
daten = [4, 2, -6, 17, 5, 12] 
print(daten[0])   # 4
print(daten[1])   # 2
...               # usw.
\end{lstlisting}

	\begin{itemize}[<+->]
		\item Welche anderen Variablen-Typen kennen Sie?
		      \begin{itemize}
			      \item Zahl: \lstinline|x=3|
			      \item Text: \lstinline|name="Hallo"|
		      \end{itemize}
		\item Wie viele Werte enthält eine Liste?
		      \begin{itemize}
			      \item Automatisch berechnet mit: \lstinline|len(daten) # 6|
			      \item Erstes Element der Liste: \lstinline|daten[0]|
			      \item Letztes Element der Liste: \lstinline|daten[len(daten)-1]|
			      \item Weshalb nicht \lstinline|daten[len(daten)]|?
		      \end{itemize}
	\end{itemize}
\end{frame}

\begin{frame}[fragile]{Listen}{Listen verändern}
	\begin{lstlisting}
# Neue Variable des Typs Liste erstellen
daten = [4, 2, -6, 17, 5, 12] 
print(daten)   # 4, 2, -6, 17, 5, 12
daten[3] = 20  # Ein Element der Liste wird verändert (welches?)
print(daten)   # ?
\end{lstlisting}
\end{frame}

\begin{frame}[fragile]{Listen}{Listen verändern}
	\begin{lstlisting}
# Neue Variable des Typs Liste erstellen
daten = [4, 2, -6, 17, 5, 12] 
print(daten)   # 4, 2, -6, 17, 5, 12
daten[3] = 20  # Ein Element der Liste wird verändert (welches?)
print(daten)   # 4, 2, -6, 20, 5, 12
\end{lstlisting}
\end{frame}


\begin{frame}[fragile]{Listen}{Index}
	\begin{lstlisting}
# Neue Variable des Typs Liste erstellen
daten = [4, 2, -6, 17, 5, 12] 
index = 0
daten[index] += 2
print(daten)
\end{lstlisting}


	\begin{itemize}[<+->]
		\item Wozu könnte die Variable \lstinline|index| \textit{dienen}? Weshalb nicht einfach nur \lstinline|daten[0]| schreiben?!?
		\item Index ist ein \textbf{Diener} und hat \textbf{keine Bedeutung}! \emoji{astonished-face}
	\end{itemize}
\end{frame}

\begin{frame}[fragile]{Durch jedes Element in einer Liste durchgehen}{Flexibler Summen-Befehl (Möglichkeit 1/2)}
	\begin{lstlisting}
		def summiere(daten):
			i = 0
			summe = 0
			n = len(daten)
			
			for _ in range(n):
				wert = daten[i]
				summe += wert
				i += 1
			
			print(summe)
			
		daten = [4, 2, -6, 17, 5, 12] 
		summiere(daten)
	\end{lstlisting}
\end{frame}


\begin{frame}[fragile]{Durch jedes Element in einer Liste durchgehen}{Flexibler Summen-Befehl (Möglichkeit 1/2, mit Kommentar)}
	\begin{lstlisting}[basicstyle=\ttfamily\scriptsize]
def summiere(daten):
	i = 0           # Hilfs-Index, um auf Elemente von daten zuzugreifen
	summe = 0       # Variable, um alle Elemente von "daten" zu summieren
	n = len(daten)  # Anzahl Elemente in "daten" (Wert?)
	
	# Durch alle Elemente von Daten durchgehen und zu Summe hinzufügen
	for _ in range(n):
		wert = daten[i] # Auf jedes Element von Daten zugreifen 
		summe += wert   # Wert zu Summe hinzufügen
		i += 1          # Hilfs-Index um 1 vergrössern (Werte?)
		
	print(summe)
	
daten = [4, 2, -6, 17, 5, 12] 
summiere(daten)
\end{lstlisting}
\end{frame}

\begin{frame}[fragile]{4 Fehler}
	Wir wollen jedes Element von \lstinline|daten| um 2 vergrössern
	\begin{lstlisting}
daten = [4, 2, -6, 17, 5, 12] 
index = 1
for _ in range(6):
	index += 1
	
print(daten)
\end{lstlisting}
\end{frame}

\begin{frame}[fragile]{4 Fehler}{Korrigiert}
	Wir wollen jedes Element von \lstinline|daten| um 2 vergrössern
	\begin{lstlisting}
daten = [4, 2, -6, 17, 5, 12] 
index = 0
for _ in range(len(daten)):
	daten[index] += 2
	index += 1
	
print(daten)
\end{lstlisting}
\end{frame}

\begin{frame}[fragile]{Skalarprodukt}
	Wo wird das Skalarprodukt verwendet?
	\begin{table}[H]
		\begin{tabular}{c||c|c}
			\textbf{Produkt} & \textbf{Menge} \only<2->{\lstinline|m|} & \textbf{Preis} \only<2->{\lstinline|p|} \\\midrule\midrule
			\emoji{banana}   & 800 g                                   & 2.- / Kg                                \\\midrule
			\emoji{apple}    & 1200 g                                  & 2.50 / Kg                               \\\midrule
			\emoji{peach}    & 2300 g                                  & 5.- / Kg                                \\
		\end{tabular}
	\end{table}
	\only<1>{Wie berechnen wir den \textbf{Gesamtpreis?}}
	\only<2>{\textbf{Skalarprodukt}: \lstinline|m[0]*p[0] + m[1]*p[1] + ... + m[-1]*p[-1]|}
\end{frame}


\begin{frame}[fragile]{Auftrag}
	\begin{itemize}
		\item \readbox{das Kapitel 6 (Datenstrukturen)}
		\item \readbox{Beispiel 6.1}
		\item \aufgabebox{Aufgaben 6.1 - 6.8}
	\end{itemize}
	\begin{table}[H]
		\centering
		\begin{tabular}{c||c|c}
			                 & [               & ]               \\\midrule\midrule
			\textbf{Windows} & \keys{\AltGr+ü} & \keys{\AltGr+!} \\\midrule
			\textbf{MacOS}   & \keys{\Alt+5}   & \keys{\Alt+6}   \\
		\end{tabular}
		\caption*{Tastatur: Eckige Klammern}
	\end{table}
\end{frame}